% ============================================================
%  SY02-46 · §3 反常层 ℒ^d_E · 印刷页 8–9（底本 p8 下段起，p9）
%  条目：Lemma 3.1、Definition 3.2
% ============================================================
\section{反常层 $\cL^d_E$}\srcpage{8}

对每个整数 $d\ge 0$，设 $\Coh^d_0(X)$ 是一般秩 $0$ 的凝聚 $\cO_X$-模 $\cM$ 的叠，也就是说 $\cM$ 具有有限支撑，且长度为 $\dim H^0(X,\cM)=d$。记 $\Coh_0(X)$ 为诸 $\Coh^d_0(X)$ 的不交并。

叠 $\Coh^d_0(X)$ 的定义对任意拟射影代数曲线 $X$ 有意义；特别地，可以取叠 $\Coh^d_0(\A^1)$。一般秩 $0$、长度为 $d$ 的凝聚 $\cO_{\A^1}$-模，等价于一个 $d$ 维向量空间连同其一个自同态。因此 $\Coh^d_0(\A^1)$ 就是 $d\times d$ 方阵的仿射空间在 $\mathrm{GL}(d)$ 共轭作用下的商叠 $[\mathfrak{gl}(d)/\mathrm{GL}(d)]$。由于任意光滑代数曲线 $X$ 在 étale 拓扑下局部同构于仿射直线 $\A^1$，$\Coh^d_0(X)$ 在 étale 拓扑下局部同构于 $[\mathfrak{gl}(d)/\mathrm{GL}(d)]$。

$X$ 上每个有效除子 $D$ 定义挠 $\cO_X$-模
\[
  \cO_{X,D} = \cO_X/\cO_X(-D),
\]
其长度等于 $D$ 的次数。并且对每个整数 $d\ge 0$，有态射
\[
  X^{(d)} \to \Coh^d_0(X),\quad D\mapsto \cO_{X,D},
\]
\srcpage{9}
它光滑，相对维数为 $d$。

任意 $\cM\in\Coh_0(X)$ 都同构于 $\bigoplus_{i\in I}\cO_{X,D_i}$，其中 $(D_i)_{i\in I}$ 是 $X$ 上有效除子的有限族。并且对每个整数 $d\ge 0$，有行列式态射
\[
  \det : \Coh^d_0(X) \to X^{(d)},
\]
它把 $\cM\cong\bigoplus_{i\in I}\cO_{X,D_i}$ 送到 $\sum_{i\in I}D_i$。记 $X^{(d)}_{\rss}$ 为 $X^{(d)}$ 中由无重除子 $D$ 构成的开集，即对任意 $x$ 有 $d_x=0$ 或 $1$；记 $\Coh^d_{0,\rss}(X)=\det^{-1}(X^{(d)}_{\rss})$ 为 $\Coh^d_0(X)$ 中对应的开集。

设 $E$ 是 $X$ 上秩 $n$ 的局部系统。关于 $E^{(d)}$ 在 $D\in X^{(d)}$ 处纤维的公式 (1.3) 表明，该纤维的秩随 $D$ 变化；但在开集 $X^{(d)}_{\rss}$ 之上，这个秩是常值，且 $E^{(d)}$ 的限制是秩 $n^d$ 的局部系统。

\begin{lemma}
对每个整数 $d\ge 0$，复形 $E^{(d)}[d]\in\Dbc(X^{(d)})$ 由可构造层 $E^{(d)}$ 放在次数 $-d$ 处构成。它是反常层，并且是其到 $X^{(d)}_{\rss}$ 上限制的中间延拓。反常层 $E^{(d)}[d]$ 在 $E$ 不可约（相应地，半单）时不可约（相应地，半单）。
\end{lemma}

把可构造层 $\det^* E^{(d)}$ 限制到 $\Coh^d_0(X)$ 的稠密开集 $\Coh^d_{0,\rss}(X)$ 上，得
\[
  \cL^d_{E,\rss} = \det^* E^{(d)}|_{\Coh^d_{0,\rss}(X)},
\]
它也是秩 $n^d$ 的局部系统，并且在 $E$ 不可约（相应地，半单）时不可约（相应地，半单）。

\begin{definition}
对每个整数 $d\ge 0$，与 $X$ 上局部系统 $E$ 相伴的反常层 $\cL^d_E$，就定义为 $\cL^d_{E,\rss}$ 到整个 $\Coh^d_0(X)$ 的中间延拓。
\end{definition}

当然，若 $E$ 不可约（相应地，半单），则 $\cL^d_E$ 不可约（相应地，半单）。此外可以验证，$\cL^d_E$ 沿态射 $X^{(d)}\to\Coh^d_0$，$D\mapsto\cO_{X,D}$ 的 $d$-平移拉回，正是反常层 $E^{(d)}[d]$。
